\section{Échelles cosmologiques associées au cycle de 108 états et à l’opérateur de Perron}
\label{rm:ch:cosmo}

\subsection{Séparation typologique des \(2\) réalisations cosmologiques}

La construction cosmologique comprend deux grandeurs internes d’origines distinctes. La première dérive de la fermeture périodique CT108 et détermine une échelle de courbure. La seconde dérive du cocycle de Perron et détermine une loi de dilution de la mémoire. Toutes deux utilisent le même tronc \(\APP\to\TRIT\to\TPK\), mais n’ont ni le même domaine ni le même sélecteur :

{\footnotesize
\begin{equation}
\TPK\longrightarrow\text{état enrichi}
\longrightarrow
\begin{cases}
\text{structure chronologique }108\longrightarrow
R_{108}\longrightarrow(H_{108},\Lambda_{108},T_{\rm dS},S_{\rm dS}),\\
\text{cocycle de Perron}\longrightarrow
(a_n,M_n)\longrightarrow(H_\varphi,s^2\propto a^{-6}).
\end{cases}
\label{rm:eq:cosmo-two-branches}
\end{equation}
}

L’égalité éventuelle des unités ou de l’ordre de grandeur n’autorise pas à identifier ces réalisations. Le premier critère de réfutation du chapitre est la prétention de les identifier sans entrelaceur préservant la structure chronologique, l’orientation, la mémoire et la carte dimensionnelle.

\subsection{Condition de périodicité du cycle de 108 états}

Soit \(\ell_0=c_{\rm int}t_0\). La période et la fréquence angulaire du retour sont

\begin{equation}
T_{108}=108t_0,
\qquad
\omega_{108}=\frac{2\pi}{108t_0},
\qquad
R_{108}=\frac{c_{\rm int}}{\omega_{108}}
=\frac{54}{\pi}\ell_0.
\label{rm:eq:cosmo-R108}
\end{equation}

La longueur du tour complet est

\begin{equation}
C_{108}=2\pi R_{108}=108\ell_0.
\label{rm:eq:cosmo-C108}
\end{equation}

Pour l’une quelconque des deux sections orientées d’action \(a\in\{\mathrm{pre},\mathrm{ret}\}\), le quantum circulaire

\[
E_{108,a}=\hbar_a\omega_{108},
\qquad
m_{108,a}=E_{108,a}/c_{\rm int}^2
\]

satisfait simultanément

\begin{equation}
\bar\lambda_C(m_{108,a})=R_{108},
\qquad
\lambda_C(m_{108,a})=C_{108}.
\label{rm:eq:cosmo-compton}
\end{equation}

La réalisation circulaire ne prend donc pas une longueur de Planck extérieure. Elle construit d’abord un rayon, puis démontre que les lectures réduite et complète de Compton en sont les deux réalisations.

\subsection{Sélecteur gravitationnel et relation Planck--TRIT}

Le transducteur gravitationnel homogène est

\begin{equation}
G_a(\theta)=\theta\frac{c_{\rm int}^5t_0^2}{\hbar_a}.
\label{rm:eq:cosmo-Gtheta}
\end{equation}

Sur le quantum circulaire, le rayon gravitationnel réduit vaut

\[
r_{{\rm g},a}(\theta)
=\frac{G_a(\theta)m_{108,a}}{c_{\rm int}^2}
=\theta\frac{\pi}{54}\ell_0.
\]

La condition \(r_{{\rm g},a}=R_{108}\) sélectionne de manière unique

\begin{equation}
\theta_{108}=\left(\frac{54}{\pi}\right)^2.
\label{rm:eq:cosmo-theta108}
\end{equation}

La coordonnée ne dépend pas du feuillet d’action : dans le produit \(G_a(\theta_{108})\hbar_a\), l’orientation avant/retour s’annule. Il en résulte

\begin{align}
\ell_P&=R_{108}=\frac{54}{\pi}\ell_0,
&
t_P&=\omega_{108}^{-1}=\frac{54}{\pi}t_0,
\label{rm:eq:cosmo-planck-length-time}\\
E_P&=\hbar_a\omega_{108}
=\frac{\pi}{54}\frac{\hbar_a}{t_0},
&
E_P&=k_B\Theta_{\rm clk}\log3.
\label{rm:eq:cosmo-planck-trit}
\end{align}

\cref{rm:eq:cosmo-planck-trit} exprime la confluence structurelle : \(\log3\) est le contenu informationnel d’une décision TRIT et, dans la carte thermique, convertit la température chronologique en quantum de Planck. L’identité n’assimile pas l’espace des états TRIT à l’espace gravitationnel ; elle démontre que deux réalisations fonctionnelles déterminent le même scalaire par des applications distinctes.

\subsection{Réalisation conditionnée de la géométrie de Sitter}

Supposons maintenant que le réalisateur de Palatini--Cartan--Holst sélectionne le rayon interne \(R_{108}\) comme rayon d’un fond de Sitter à quatre dimensions. L’hypothèse est géométrique et demeure explicite : elle ne se déduit pas d’une coïncidence numérique avec \(H_0\). Sous cette hypothèse,

\begin{equation}
H_{108}=\frac{c_{\rm int}}{R_{108}}
=\frac{\pi}{54t_0},
\qquad
\Lambda_{108}=\frac{3}{R_{108}^2}
=\frac{\pi^2}{972\ell_0^2}.
\label{rm:eq:cosmo-ds-H-Lambda}
\end{equation}

La température de Gibbons--Hawking et l’entropie d’horizon deviennent

\begin{align}
k_BT_{\rm dS}
&=\frac{\hbar_a H_{108}}{2\pi}
=\frac{\hbar_a}{108t_0},
&
\frac{T_{\rm dS}}{\Theta_{\rm clk}}
&=\frac{\log3}{2\pi},
\label{rm:eq:cosmo-ds-temperature}\\
\frac{S_{\rm dS}}{k_B}
&=\frac{\pi R_{108}^2}{\ell_P^2}=\pi.
\label{rm:eq:cosmo-ds-entropy}
\end{align}

\(\pi\) de \cref{rm:eq:cosmo-ds-entropy} n’est pas introduit comme valeur d’une entropie à reproduire : il apparaît parce que le sélecteur circulaire impose \(R_{108}=\ell_P\) et que l’aire de la sphère d’horizon contient le facteur géométrique \(4\pi\).

\begin{proposition}[Famille sans dimension CT108--de Sitter]
Sous la réalisation de Sitter déclarée, le quadruplet
\[
\left(
H_{108}t_0,
\Lambda_{108}\ell_0^2,
\frac{T_{\rm dS}}{\Theta_{\rm clk}},
\frac{S_{\rm dS}}{k_B}
\right)
=\left(
\frac\pi{54},
\frac{\pi^2}{972},
\frac{\log3}{2\pi},
\pi
\right)
\]
est une sortie sans dimension exacte.
\end{proposition}

\begin{falsifier}
La réalisation de Sitter est invalidée si le fond sélectionné ne satisfait pas \(R_{\mu\nu}=\Lambda_{108}g_{\mu\nu}\), si le réalisateur PCH détermine un rayon différent de \(R_{108}\), ou si les cartes utilisées dans \cref{rm:eq:cosmo-ds-H-Lambda,rm:eq:cosmo-ds-temperature} ne partagent pas \((c_{\rm int},t_0,\hbar_a)\). Aucun de ces échecs ne modifie le théorème circulaire \cref{rm:eq:cosmo-R108,rm:eq:cosmo-theta108}.
\end{falsifier}
